\documentclass[11pt]{article}
\usepackage[margin=1in]{geometry}
\usepackage{charter}
\usepackage[hidelinks]{hyperref}
\usepackage{setspace}

\pagestyle{empty}

\begin{document}

\noindent
\textbf{Vijayakumar Thangavel Mahalingam} \\
Department of Pharmacy Practice \\
SRM College of Pharmacy, Faculty of Medicine and Health Science \\
SRM Institute of Science and Technology \\
Kattankulathur, Chennai, India \\
ORCID: \href{https://orcid.org/0000-0002-7749-3224}{0000-0002-7749-3224} \\
Email: \href{mailto:vijayakm2@srmist.edu.in}{vijayakm2@srmist.edu.in}

\bigskip

\noindent
July 3, 2026

\bigskip

\noindent
Dr.\ Guang Hu \\
Handling Editor, \textit{Scientific Reports} \\
Nature Portfolio

\bigskip

\noindent
Dear Dr.\ Hu,

\onehalfspacing

We are grateful for the opportunity to submit a revised version of our manuscript, now titled \textbf{``Separating effect size from statistical evidence in network-proximity rankings under target-count asymmetry: a controlled liver-interactome audit''} (formerly submitted under a metric-correction framing), for consideration as a Research Article in \textit{Scientific Reports}.

We thank you and both reviewers for an exceptionally constructive critique. The reviews prompted us to reconsider the central framing of the work, and the revision is substantially stronger as a result. A point-by-point response accompanies this letter; the principal changes are summarised below.

\bigskip
\noindent \textbf{Reframing the central claim.}
Reviewer~2 noted that the dependence of proximity Z-scores on target-set size reflects expected statistical precision rather than a flaw in the metric. We have revised the manuscript to make this distinction explicit throughout. The revised thesis is interpretive and, we believe, more useful to the community: network-proximity Z-scores are valid \emph{evidence} statistics whose \emph{magnitude} should not be read as a cross-compound \emph{effect-size} ranking when target counts differ sharply. This is the effect-size-versus-significance distinction of the ASA statement on p-values (Wasserstein \& Lazar, 2016), transposed to network medicine, and it is fully consistent with Guney et al.\ (2016) and Menche et al.\ (2015). Language describing variance shrinkage as a general ``bias'' or ``artifact'' has been removed.

\bigskip
\noindent \textbf{New analyses and a more conservative interpretation of the headline result.}
We formalise the null-variance shrinkage as a scaling result observed across the evaluated proximity and influence metrics, add the requested parameter-sensitivity analyses, and revalidate our proximity implementation against Guney's published toolbox code (\texttt{emreg00/\allowbreak toolbox}), reproducing $d_c$ exactly. We also add a new analysis (\S2.5, Table~3, Figure~6) separating \emph{direct} target--disease overlap from \emph{propagated} network influence. This shows that much of the raw per-target advantage reflects direct overlap, while the propagated advantage is more modest ($\sim$1.5$\times$). We report these components separately and have revised the interpretation accordingly.

\bigskip
\noindent \textbf{Code and data deposition (per your request).}
The complete pipeline, curated data, pinned dependency versions, input checksums, and figure scripts are available on GitHub and have been deposited in Zenodo under the DOI-linked record \url{https://doi.org/10.5281/zenodo.21166752}. The manuscript now contains a dedicated \textbf{Code Availability} section. The principal quantitative results are regenerated from the committed data by verification scripts included in the repository.

\bigskip
\noindent \textbf{Presentation.}
The manuscript has been rebuilt to journal standard: decimal-aligned tables, consistent caption and float styling, line numbers for review, and figures regenerated without the interpretive overlay text the reviewers flagged.

\bigskip
\noindent \textbf{Change to the author list.}
We respectfully wish to inform you of a correction and expansion of the author list relative to the originally submitted manuscript, and we kindly request the editor's approval. The original manuscript was submitted under a single author (Antony Bevan). This inadvertently omitted Vijayakumar Thangavel Mahalingam, who mentored and supervised the study from its inception; the omission arose from the first author's unfamiliarity with authorship and submission conventions at the time of the original submission, for which we sincerely apologise. Separately, Sarputheen~M contributed to the present revision --- specifically the STRING text-mining sensitivity analysis and the direct-versus-propagated influence decomposition, together with critical revision of the manuscript --- and is included in recognition of that work. Accordingly, Vijayakumar Thangavel Mahalingam now serves as corresponding author. All three authors have read and approved the revised manuscript and consent to the revised author list, author order, and change of corresponding author; Antony Bevan specifically consents to no longer serving as corresponding author. We would be glad to complete any authorship-change declaration and to provide any further documentation the editorial office may require.

\bigskip
\noindent \textbf{Suggested reviewers:}
\begin{itemize}
    \item \textbf{Feixiong Cheng}, Cleveland Clinic (Genomic Medicine Institute) -- chengf@ccf.org \\ \textit{Expert in network-based drug repurposing and interactome proximity.}
    \item \textbf{Minjun Chen}, National Center for Toxicological Research, FDA -- minjun.chen@fda.hhs.gov \\ \textit{Creator of DILIrank; expert in predictive liver toxicology.}
    \item \textbf{Patrick Aloy}, IRB Barcelona -- patrick.aloy@irbbarcelona.org \\ \textit{Expert in network biology and computational systems pharmacology.}
\end{itemize}

\bigskip
\noindent
We confirm that this manuscript has not been published elsewhere and is not under consideration by another journal. We have had no prior discussions with a \textit{Scientific Reports} Editorial Board Member about the work described in the manuscript. We have endeavoured to address each reviewer and editorial point with care and to keep this letter brief in respect of your time. We are sincerely grateful for the consideration you and the reviewers have given our work; we hope the revision does justice to that effort. We remain fully open to any further revisions, analyses, or documentation that you or the editorial office may require, and we will respond promptly to any additional requests.

\bigskip
\noindent
Sincerely,

\vspace{1em}

\noindent
\textbf{Vijayakumar Thangavel Mahalingam} \\
Corresponding Author

\end{document}
